% Three guards that could not fire. Sits inside §3 (sec:calibration) with no
% \section of its own — see PATCHES.md E1/G4.

The two defective validation statistics --- the recovery curve of
\S\ref{sec:blind} and the grid above --- are not isolated. Our working rule holds that \textbf{a check unable to fail is worse
than no check}, because such a check turns an absence of evidence into a green
light. Three more guards shipped and got withdrawn, and all three survive
correct code. The third we wrote ourselves, to this rule, after the other two
were already written up. A fifth statistic is not a guard at all --- it is the
interval every one of these readings ends in, and it closes this section.
Five others were plain implementation errors: three sit in
Appendix~\ref{app:withdrawn} as record rather than finding, and two are named
there and not described.

\paragraph{A confound diagnostic unable to fire, and a repair unable to fire
either.} Our labelling axis defines the cell state by the \emph{absence} of
markers, because every positive marker doubles as a measurement target and a
leakage invariant bars a target from serving as a label. Absence tracks
sequencing depth closely, so the obvious guard correlates the maturity call
against depth and flags anything above $|\rho| = 0.20$. The correlation between
a continuous variable and a \emph{binary} label of prevalence $p$ carries a
bound of $\sqrt{3p(1-p)}$. Below $p = 1.3516\%$ the 0.20 tolerance sits above
the bound and the diagnostic reads clean whatever the data does. Our rarest rung
held 8 positive cells in 5{,}564, a bound of 0.066, and got filed as cleanest of
four.

The same collapse happens at the other end. On the coarsest resolution every
epithelial cell is called mature, prevalence is exactly $1.000$, and the bound
is exactly zero: the correlation is undefined on all 64 patient-by-axis rows and
the check reports nothing at all. \textbf{Degenerate labels are invisible to
this check at either extreme}, and the two extremes are the coarsest and the
finest resolution we ran.

So the raw statistic does not rank. A resolution whose ceiling is 0.066 and one
whose ceiling is 0.866 are not comparable on $|\rho|$ at all --- the same
observed correlation carries an order of magnitude more evidence at one than at
the other, and the tolerance is a single number applied to both.
Cleanest-of-four was a ranking of ceilings.

Our repair inherited the defect. The repaired diagnostic reported the attainable
bound, but took the worst correlation and the largest bound as
\emph{independent} maxima over the two arms, so the statistic and its bound
could come from different arms --- and it read clean when they did. Correcting
the pairing took the reachability flag with it, and at the rarest rung that flag
had called the test capable of firing on \textbf{64 of 64} combinations where 16
is correct.
\textbf{The blind spot sits exactly where the science is.} Rare populations are
the ones worth resolving, and the ones this check never flags.

\paragraph{A pre-registered criterion unable to pass when the hypothesis holds.}
We built one of four gate criteria so the condition it tested gets violated
\emph{by} the effect the study exists to detect, under every population
definition anyone tried. Its pre-committed consequence never fired, because
acting on a criterion proven unable to pass would have been caused by the
pre-registration rather than prevented by it. Pre-registration protects you
against choosing the analysis after seeing the result. It offers no protection
against an unsatisfiable criterion, and standard practice never asks.


\paragraph{A premise check whose sensitivity was a property of the gene it
controlled.} Forbidden a mechanism claim from the decomposition, we asked the
question a second way: in patients whose mature cells survive in \emph{both}
arms, does the marker fall within those cells? That reading is licensed only if
the two arms' surviving mature cells are comparable, so we gated it on a premise
check --- housekeeping controls compared between arms against a tolerance, run
before anything else is read. On detection rate the controls pass. Across three
cohorts and 42 patients, \textbf{83 of 84} patient-gene control rows sit inside a
$0.10$ tolerance, and the premise reads \emph{holds}.

The controls had moved. Those same genes shift by factors of $1.19$ to $1.64$ in
per-cell output between the arms, and detection rate cannot see it, because at a
reference detection rate of $0.967$ to $1.000$ it has nowhere to go. Detection
under a Poisson model is $d' = 1 - (1-d)^{s}$ for a fold change $s$, so a gene
detected in 99\% of cells whose output \emph{halves} still reads $0.90$ --- inside
the tolerance by construction. \textbf{The tolerance again sat above the
statistic's attainable range}: this is the confound diagnostic above, in a guard
the same people wrote four days after writing that paragraph, to the rule stated
at the top of this section. Reassessed on log$_2$ expression, the primary
cohort refuses the premise it had just granted.

The refusal did not hold either. The verdict was a point estimate against a
fixed line, so it inherited exactly the defect the estimator is gated against:
one cohort's control read $+0.486$ under one seed and $+0.540$ under another,
against a tolerance of $0.5$ --- the same data, a different draw, opposite
verdicts. Bootstrapped over patients, every cohort's control interval straddles
the tolerance: $-0.712\,[-1.250, -0.289]$, $+0.540\,[+0.012, +1.112]$,
$+0.431\,[+0.216, +0.640]$. So the premise check ended with the shape of the
estimator it exists to gate --- satisfied, refused, and \textbf{undecided} --- and
returns undecided on all three cohorts. The mechanism question is unresolved on
this data. \textbf{That is the finding a two-state check would not have let us
reach}: it said \emph{holds}, and we would have published a mechanism.

\paragraph{And the fifth is not a guard at all. It is the interval.}
\label{sec:interval}
Every reading in this paper ends in a percentile bootstrap over patients. The
bootstrap distribution of a mean has standard deviation $s\sqrt{(n-1)/n}$ --- the
plug-in one --- so the percentile interval is approximately
$\bar{x} \pm z\,s\sqrt{(n-1)/n}/\sqrt{n}$ against the correct
$\bar{x} \pm t_{n-1}\,s/\sqrt{n}$, and the ratio of the two widths is
\begin{center}
$z\sqrt{(n-1)/n} \,/\, t_{n-1}$ \quad --- \quad \textbf{a function of $n$ alone,
with no data in it.}
\end{center}
Two errors running the same way: a normal quantile where a $t$ quantile is
needed, and a biased standard deviation where an unbiased one is needed. At
$n = 44$ the interval is $0.96\times$ the width it claims and a nominal 5\% test
rejects 5.9\% of the time; at $n = 20$, $0.91\times$ and 7.1\%; at $n = 10$,
$0.82\times$ and 9.6\%; at $n = 4$, $0.53\times$ and 18.8\%. Simulation at 1{,}500
trials per cell agrees with the closed form to within 1.5 percentage points at
every $n$, and the excess is positive everywhere --- the gap is the skew a normal
approximation drops, so the closed form is a floor.

\textbf{Neither sophisticated repair works, and one is worse.} BCa is
miscalibrated in 17 of 20 cells against the percentile's 14, because its bias
correction and jackknife acceleration are estimated from the same ten numbers.
Raising the bootstrap replicate count does nothing at all: the error is not
Monte Carlo error, so more replicates estimate the wrong thing more precisely.
The plain Student-$t$ interval is calibrated in 19 of 20 cells and conservative
in the twentieth, worst case 5.9\%.

This is the same failure as the four above, one layer out. A guard that cannot
fail turns absence of evidence into a green light; an interval narrower than it
claims turns noise into a finding. Neither announces itself, and both are
functions of something other than the quantity a reader thinks they measure ---
\textbf{which is why the diagnostic this paper proposes is a question about the
statistic rather than a closer look at its output.} We did not find this by
inspecting intervals. We found it by asking what ours was a function of.
